\appendix
\section{Mathematical Derivations}
\label{chapter:A}

% block-id: A.report_derivation.heading
\hypertarget{A.report_derivation.heading}{}
\subsection{Analytical Normalization, Circuit Constraints, and Element Integrals}

% block-id: A.report_derivation.p1
\hypertarget{A.report_derivation.p1}{}
Round-wire normalization first requires 2π\(\int\)\(_0\)\(^a\) r Jz(r) dr=It. The identity d[r I\(_1\)(κr)]/dr=κr I\(_0\)(κr) then gives the normalization denominator. Dividing the boundary axial electric field Jz(a)/σ by It yields the internal impedance per unit length. The independent Joule integral is compared with |It|² Re(zint) under the RMS convention. High-precision quadrature error checks only the stated analytical formula; finite ends, external return inductance, temperature changes, and insulation are not included in that analytical accuracy.

% block-id: A.report_derivation.p2
\hypertarget{A.report_derivation.p2}{}
Linear-triangle recovery uses the actual physical area aT. With nodal eₖ=v−uₖ, complex element current is σaT\(\sum\)eₖ/3 and Joule power is σaT[\(\sum\)|eₖ|²+|\(\sum\)eₖ|²]/12. The latter follows from the consistent mass matrix with diagonal coefficient two and off-diagonal coefficient one. Original coordinates are scaled by 0.001 m, so area carries the square of that scale; conversion to millimetres for display does not change integration. Material markers distinguish individual strands, return, and air, and the return uses its own v.

% block-id: A.report_derivation.p3
\hypertarget{A.report_derivation.p3}{}
Necessity and sufficiency of the common-terminal criterion depend on the same linear system and uniqueness of its solution. Necessity follows by substituting equal currents into the voltage equation. Sufficiency follows because a vector removed by Q can be represented as V1 and then satisfies the augmented system together with total current. Conductor contacts, additional terminal impedances, nonlinear material, or a missing induced term change that system. Reuse of the criterion in a new model therefore begins by reconstructing Z, e, and the constraints, rather than copying an earlier conclusion.

% block-id: A.sharing_derivation.heading
\hypertarget{A.sharing_derivation.heading}{}
\subsection{Derivation of the uniform complex-current condition}

% block-id: A.sharing_derivation.p1
\hypertarget{A.sharing_derivation.p1}{}
If a uniform-current solution exists, substituting i=(It/N)1 into the port equation makes all components of (It/N)Z1+e equal, establishing necessity of the projection condition. Conversely, when the projection condition holds, this vector equals a scalar V times 1. The uniform-current vector then satisfies both the port equation and the total-current constraint; uniqueness of the augmented solution makes it the actual solution, establishing sufficiency. Numerical results use sqrt(mean(|ik−It/N|²))/|It/N| for normalized RMS complex-current deviation and show individual magnitudes and phases.

% block-id: A.method.sources
\hypertarget{A.method.sources}{}
The modified Bessel functions used for the round-wire reference follow the definitions in \cite{BESSEL}. Partial-inductance construction follows the PEEC framework \cite{RUEHLI1974}; the face-flux approximation is related to the Raviart–Thomas family \cite{RT1977}. GMRES supplies the iterative linear-solver framework \cite{GMRES1986}, while the finite-element implementation and boundary treatment are documented in \cite{FEMDOC}.
